\documentclass{article}
\usepackage{mathtools} 
\usepackage{fontspec}
\usepackage[UTF8]{ctex}
\usepackage{amsthm}
\usepackage{mdframed}
\usepackage{xcolor}
\usepackage{amssymb}
\usepackage{amsmath}


% 定义新的带灰色背景的说明环境 zremark
\newmdtheoremenv[
  backgroundcolor=gray!10,
  % 边框与背景一致，边框线会消失
  linecolor=gray!10
]{zremark}{说明}

% 通用矩阵命令: \flexmatrix{矩阵名}{元素符号}{行数}{列数}
\newcommand{\flexmatrix}[4]{
  \[
  #1 = \begin{pmatrix}
    #2_{11}     & #2_{12}     & \cdots & #2_{1#4}   \\
    #2_{21}     & #2_{22}     & \cdots & #2_{2#4}   \\
    \vdots      & \vdots      & \ddots & \vdots     \\
    #2_{#31}    & #2_{#32}    & \cdots & #2_{#3#4}
  \end{pmatrix}
  \]
}

% 简化版命令（默认矩阵名为A，元素符号为a）: \quickmatrix{行数}{列数}
\newcommand{\quickmatrix}[2]{\flexmatrix{A}{a}{#1}{#2}}

\begin{document}
\title{注释 14.2}
\author{张志聪}
\maketitle

\begin{zremark}
  证明：定理14.11
\end{zremark}
\textbf{证明：}

设泰勒级数的和函数为$S(x)$，部分和函数为
\begin{align*}
  S_n(x) = \sum \limits_{k = 0}^n \frac{f^{(k)}(x_0)}{k!}(x - x_0)^k
\end{align*}
于是泰勒公式余项可表示为：
\begin{align*}
  R_n(x) = f(x) - S_n(x)
\end{align*}
注：这里的泰勒公式的余项是第六章定理6.9下定义的，
不是函数项级数的余项。

\begin{itemize}
  \item 必要性

        已知$f(x)$在区间$(x_0 - r, x_0 + r)$上等于它的泰勒级数的和函数，
        即$S(x) = f(x)$，
        由和函数的定义可知，对任意$x \in (x_0 - r, x_0 + r)$，有
        \begin{align*}
          \lim\limits_{n\to\infty} S_n(x) = f(x)
        \end{align*}
        从而
        \begin{align*}
          \lim\limits_{n\to\infty} R_n(x)
           & = \lim\limits_{n\to\infty} f(x) - S_n(x) \\
           & = f(x) - f(x)                            \\
           & = 0
        \end{align*}

  \item 充分性

        对任意$x \in (x_0 - r, x_0 + r)$，有
        \begin{align*}
          \lim\limits_{n\to\infty} R_n(x) = 0
        \end{align*}
        于是
        \begin{align*}
          S(x)
           & = \lim\limits_{n\to\infty} S_n(x)                 \\
           & = \lim\limits_{n\to\infty} f(x) - (f(x) - S_n(x)) \\
           & = \lim\limits_{n\to\infty} f(x) - R_n(x)          \\
           & = f(x) - 0                                        \\
           & = f(x)
        \end{align*}
\end{itemize}

\begin{zremark}
  函数$f$的泰勒展开式的唯一性。
\end{zremark}
\textbf{证明：}
考虑$x_0 = 0$时，设$f$的泰勒展开式为
\begin{align*}
  f(x) = \sum \limits_{n = 0}^{\infty} a_n x^n
\end{align*}
由于$\sum \limits_{n = 0}^{\infty} a_n x^n$也是幂级数，
由定理14.8推论2可知，
\begin{align*}
  a_0 = f(0)              \\
  a_1 = f'(0)             \\
  a_2 = \frac{f''(0)}{2!} \\
  \vdots                  \\
  a_n = \frac{f^{(n)}(0)}{n!}
\end{align*}
因此，系数$a_n$被唯一确定，即泰勒展开式是唯一的。

$x_0 \neq 0$时，令$y = x - x_0$，
于是，$f$的泰勒展开式为
\begin{align*}
  f(x) = \sum \limits_{n = 0}^{\infty} a_n y^n
\end{align*}
由之前的讨论可知，其是唯一的。

% \begin{zremark}
%   例8的疑问

%   已知$\ln(1 + x)$的幂级数展开式为$\sum \limits_{n = 1}^{\infty} (-1)^{n - 1} \frac{x^n}{n}, x \in (-1, 1]$，
%   为什么通过变量替换计算$(1 - x) \ln (1 - x)$的幂级数展开式后，
%   教材讨论$x$的取值范围时，为什么没有采用直接通过$-x \in (-1, 1]$计算出$x \in [-1, 1)$.
% \end{zremark}
% \textbf{证明：}
% 这种替换可能会影响端点处的敛散性，要验证端点处的情况。

\end{document}